% One scheduler step of the capstone engine, worked by hand. block_size 4, token budget 6.
% Request A (prompt 7 tokens, 3 generated) processes its 3rd generated token at position 9; request B
% (prompt 9 tokens) has 4 prompt tokens cached and receives a prefill chunk of 5.
% Illustrative values, computed with the engine's rules.
\begin{tikzpicture}[
  cell/.style={draw=BookBlue, minimum width=6.2mm, minimum height=6.2mm, inner sep=0pt, font=\scriptsize\ttfamily},
  cached/.style={cell, fill=BookPale},
  new/.style={cell, fill=BookGold!70},
  empty/.style={cell, fill=white, draw=black!25},
  blk/.style={font=\scriptsize\sffamily, text=BookBlue},
  lab/.style={font=\small\sffamily, anchor=west},
  num/.style={circle, fill=BookBlue, text=white, inner sep=1pt, font=\scriptsize\bfseries, minimum size=4mm},
]
% ---- request A: positions 0..9, blocks 3, 0, 5 (block table [3, 0, 5])
\node[lab] at (-0.3, 3.9) {\textbf{A}: prompt 7 tokens, 3 generated; table [3, 0, 5]};
\foreach \p in {0,...,8} {\node[cached] at ({0.3 + 0.62*\p + 0.25*floor(\p/4)}, 3.2) {\p};}
\node[new] at ({0.3 + 0.62*9 + 0.25*2}, 3.2) {9};
\foreach \p in {10,11} {\node[empty] at ({0.3 + 0.62*\p + 0.25*2}, 3.2) {};}
\node[blk] at (1.23, 2.7) {block 3}; \node[blk] at (3.96, 2.7) {block 0}; \node[blk] at (6.69, 2.7) {block 5};
% ---- request B: positions 0..8 of prompt, blocks 1, 4, 2
\node[lab] at (-0.3, 1.9) {\textbf{B}: prompt 9 tokens, 4 cached; table [1, 4, 2]};
\foreach \p in {0,...,3} {\node[cached] at ({0.3 + 0.62*\p + 0.25*floor(\p/4)}, 1.2) {\p};}
\foreach \p in {4,...,8} {\node[new] at ({0.3 + 0.62*\p + 0.25*floor(\p/4)}, 1.2) {\p};}
\foreach \p in {9,10,11} {\node[empty] at ({0.3 + 0.62*\p + 0.25*2}, 1.2) {};}
\node[blk] at (1.23, 0.7) {block 1}; \node[blk] at (3.96, 0.7) {block 4}; \node[blk] at (6.69, 0.7) {block 2};
% ---- legend
\node[cached, minimum width=4mm, minimum height=4mm] at (0.3, -0.35) {}; \node[lab, font=\scriptsize\sffamily] at (0.5, -0.35) {already in the cache};
\node[new, minimum width=4mm, minimum height=4mm] at (3.7, -0.35) {}; \node[lab, font=\scriptsize\sffamily] at (3.9, -0.35) {computed this step};
\node[empty, minimum width=4mm, minimum height=4mm] at (6.8, -0.35) {}; \node[lab, font=\scriptsize\sffamily] at (7.0, -0.35) {reserved at admission};
% ---- the plan, step by step
\begin{scope}[xshift=9.3cm]
\node[lab] at (0, 3.9) {\textbf{Plan} (token budget 6)};
\node[num] at (0.2, 3.25) {1}; \node[lab, font=\small\sffamily, align=left] at (0.5, 3.25) {decode first: A gets 1 token\\at position 9 \quad budget 6 $\to$ 5};
\node[num] at (0.2, 2.25) {2}; \node[lab, font=\small\sffamily, align=left] at (0.5, 2.25) {then prompts: B gets its\\remaining 5 tokens \quad 5 $\to$ 0};
\node[num] at (0.2, 1.25) {3}; \node[lab, font=\small\sffamily, align=left] at (0.5, 1.25) {one forward pass, $n = 6$ tokens;\\A's and B's last tokens want logits};
\node[num] at (0.2, 0.25) {4}; \node[lab, font=\small\sffamily, align=left] at (0.5, 0.25) {A samples its 4th token;\\B samples its 1st: prefill done};
\end{scope}
\end{tikzpicture}
